%% kapa: exam=2 round=26 year=2015 period=3 subject=compensation_law subject_ko=감정평가 및 보상법규 seq=4 number=4 points=10
%% source: Q-Net 감정평가사 기출문제 (cst00309) artlSeq=5149601 file=extracted/1008/5149601/2015년 제26회 감정평가사 2차시험 기출문제/3교시(감정평가 및 보상법규).pdf
%% conversion: kordoc 4.24.51 -> tools/gen2.py (auto); review=proofread-gpt-sol; flags=none
\begin{longquestion}{【문제 4】}
감정평가사 甲은 토지소유자 乙로부터 그 소유의 토지(이하 ‘이 사건 토지’라고 한다)를 물류단지로 조성한 후에 형성될 이 사건 토지에 대한 추정 시가를 평가하여 달라는 감정평가를 의뢰받아 1천억 원으로 평가하였다(이하 ‘이 사건 감정평가’라고 한다). 甲은 그 근거로 단순히 인근 공업단지 시세라고 하며 공업용지 평당 3백만 원 이상이라고만 감정평가서에 기재하였다. 그러나 얼마 후 이 사건 토지에 대한 경매절차에서 법원의 의뢰를 받은 감정평가사 丙은 이 사건 토지의 가격을 1백억 원으로 평가하였다. 평가금액 간에 10배에 이르는 현저한 차이가 발생하자 사회적으로 문제가 되었다. 이에 국토교통부장관은 적법한 절차를 거쳐 甲에게 “부동산의 적정한 가격을 산정하기 위해서는 정확한 자료를 검토하고 이를 기반으로 가격형성요인을 분석하여야 함에도 그리하지 않은 잘못이 있다.”는 이유로 징계를 통보하였다. 이에 대해 甲은 이 사건 감정평가는 미래가격 감정평가로서 비교표준지를 설정할 수 없어 부득이하게 인근 공업단지의 시세를 토대로 평가하였던 것이고, 미래가격 감정평가에는 구체적인 기준이 따로 없으므로 일반적인 평가방법을 따르지 않았다고 해서 자신이 잘못한 것은 아니라고 주장한다. 甲의 주장은 타당한가? (10점) \par
\end{longquestion}
